\begin{textsample}{POS Dim 1 – 2017\_c5 – Score 15.00 – 2017\_c5/enade\_2017\_geo\_3150a0cc00.txt}  \label{ex:f1_pos_058}
CAPÍTULO 5 DISTRIBUIÇÃO DOS CONCEITOS 5.1 \textbf{PANORAMA} NACIONAL DA DISTRIBUIÇÃO DOS CONCEITOS O Conceito Enade é um indicador de qualidade que avalia os cursos por \textbf{intermédio} dos desempenhos dos estudantes a partir dos resultados do Enade.
Os valores possíveis vão de 1 ( pior situação ) a 5 ( melhor situação ).
Os cursos com apenas um ou sem nenhum concluinte participante não obtêm o Conceito Enade, ficando Sem Conceito ( SC ).
A Tabela 5.1a apresenta a quantidade e distribuição de cursos de Licenciatura de Geografia participantes do Enade/2017, por faixa de conceito e Grande Região.
Enfatiza -se, mais uma vez, que a diferença entre os cursos \textbf{tabulados} neste capítulo e no capítulo 2 corresponde aos cursos Sem Conceito ( SC ), ou seja àqueles, em princípio, sem alunos concluintes que houvessem participado da prova ou àqueles que tivessem tão somente um aluno concluinte¹⁸.
Observando -se os dados da Tabela 5.1a, nota -se que, dos 241 cursos participantes, 88 ( 36,5 \% ) classificaram -se com conceito 3, o valor modal.
Esse foi também o conceito modal em três das regiões : Norte ( 50,0 \% ), Nordeste ( 34,9 \% ) e Sudeste ( 35,7 \% ).
Na região Sul a moda foi o conceito 4, atribuído a 39,5 \% dos cursos da região.
Já na região Centro-Oeste, a moda foi o conceito 2, atribuído a 43,8 \% dos cursos da região.
O conceito 4 foi o segundo mais frequente em nível nacional ( 27,4 \%, correspondendo a 66 cursos ) e o conceito 2, o terceiro ( 22,8 \%, correspondendo a 55 cursos ).
Houve, ainda, 17 cursos ( 7,1 \% ) que receberam conceito 5 e nove cursos ( 3,7 \% ) que receberam conceito 1.
Dos 241 cursos de Geografia ( Licenciatura ), seis ( 2,5 \% ) ficaram Sem Conceito ( SC ). 18 Estes cursos com somente um concluinte não foram considerados no capítulo 2.
A região Norte participou com 24 cursos ou 10,0 \% do total nacional.
Desses, 12 cursos receberam o conceito 3, o que \textbf{equivale} a 50,0 \% do total regional, além de ser o conceito modal na região, como já comentado.
A dois cursos ( 8,3 \% ) atribuiu -se conceito 1, a sete cursos ( 29,2 \% ), conceito 2 e a dois cursos ( 8,3 \% ), conceito 4.
Um curso ficou Sem Conceito ( SC ) e nenhum curso recebeu conceito 5 nessa região.
A região Nordeste participou com 63 cursos ou 26,1 \% do total nacional.
Nessa região, 22 cursos ( 34,9 \% em termos regionais ) obtiveram conceito 3, o conceito modal para a região, como já comentado.
Os conceitos 1 e 2 foram atribuídos, respectivamente, a quatro e a 19 cursos ( 6,3 \% e 30,2 \%, respectivamente ).
Ainda nessa região, 16 cursos ( 25,4 \% ) foram avaliados com o conceito 4, e dois cursos ( 3,2 \% ), com conceito 5.
Nenhum curso ficou Sem Conceito ( SC ).
Dos 84 cursos participantes da região Sudeste ( 34,9 \% do total nacional ), o conceito 3 foi o conceito modal, como já anteriormente assinalado, tendo sido obtido por 30 cursos ( 35,7 \% do total regional ).
O conceito 1 foi atribuído a dois cursos ( 2,4 \% ) e o conceito 2, a 12 cursos ( 14,3 \% ).
Receberam os conceitos 4 e 5, respectivamente, 27 cursos ( 32,1 \% ) e dez cursos ( 11,9 \% ).
Nessa região, três cursos ( 3,6 \% ) ficaram Sem Conceito ( SC ).
Os 38 cursos da região Sul corresponderam a 15,8 \% do total nacional.
Como já \textbf{indicado}, houve predominância do conceito 4, atribuído a 39,5 \% dos cursos da região ( conceito modal ), ou seja, atribuído a 15 dos 38 cursos participantes na região Sul.
O conceito 1 não foi atribuído a nenhum curso, o conceito 2, a três cursos ( 7,9 \% ), o conceito 3, a 14 cursos ( 36,8 \% ) e o conceito 5, a quatro cursos ( 10,5 \% ).
Dois cursos ( 5,3 \% ) ficaram Sem Conceito ( SC ) na região Sul.
Já dos 32 cursos participantes na região Centro-Oeste ( 13,3 \% do total nacional ), 14 receberam conceito 2, o conceito modal ( 43,8 \% ), como já fora destacado.
A um curso foi atribuído conceito 1 ( 3,1 \% ), e a dez cursos ( 31,3 \% ), conceito 3.
Os conceitos \textbf{restantes}, 4 e 5, foram atribuídos, respectivamente, a seis e um curso ( 18,8 \% e 3,1 \% ).
Nessa região, nenhum curso ficou Sem Conceito ( SC ).
O Gráfico 5.1a apresenta a distribuição \textbf{cumulativa} do Conceito Enade dos cursos de Licenciatura segundo Grande Região.
Num gráfico de distribuição \textbf{cumulativa}, em cada valor de \textbf{abscissa} ( eixo x ), o valor apresentado no eixo das \textbf{ordenadas} ( eixo y ) é o valor acumulado da \textbf{variável}.
Por exemplo, a região Norte ( linha azul ) apresenta um pouco mais de 90 \% dos cursos com conceito 3 ou abaixo deste valor.
A região Nordeste ( linha verde ), por outro lado, apresenta um pouco mais de 70 \% dos cursos na mesma situação : com conceito 3 ou abaixo.
Na região Norte ( linha azul ) todos os cursos têm conceitos 4 ou abaixo disso, já que a \textbf{poligonal} \textbf{cumulativa} alcança 100 \% para o conceito 4.
Em linhas gerais, num gráfico de distribuição \textbf{cumulativa}, \textbf{poligonais} mais à esquerda representam Grandes Regiões com pior distribuição, e \textbf{poligonais} mais à direita, Grandes Regiões com melhores conceitos.
A região Norte ( linha azul ) e a região Centro-Oeste ( linha preta ) apresentam os cursos com conceitos mais baixos, ao passo que as regiões Sudeste ( linha vermelha ) e Sul ( linha \textbf{roxa} ) apresentam as distribuições com valores maiores.
Cursos em IES situadas nas regiões Nordeste ( linha verde ) estão em uma situação intermediária.
A Tabela 5.1b apresenta a quantidade e distribuição de cursos de Bacharelado de Geografia participantes do Enade/2017, por faixa de conceito e Grande Região.
Enfatiza -se, mais uma vez, que a diferença entre os cursos \textbf{tabulados} neste capítulo e no capítulo 2 corresponde aos cursos sem conceito, ou seja àqueles, em princípio, sem alunos concluintes que houvessem participado da prova ou àqueles que tivessem tão somente um aluno concluinte19.
Observando -se os dados da Tabela 5.1b, nota -se que, dos 78 cursos participantes, 30 ( 38,5 \% ) classificaram -se com conceito 4, o valor modal.
Esse foi também o conceito modal nas regiões Sudeste ( 51,9 \% ) e Sul ( 41,2 \% ).
Na região Norte, a moda ficou distribuída em dois valores contíguos : conceitos 2 e 3, cada um com 33,3 \% dos cursos da região ( três cursos cada conceito ).
Na região Centro-Oeste, a moda também ficou distribuída em dois valores contíguos : conceitos 3 e 4, cada um também com 33,3 \% dos cursos da região ( três cursos cada conceito ).
Na região Nordeste, a moda foi o conceito 3 com 43,8 \% dos cursos da região, equivalente a 7 cursos.
O conceito 3 foi o segundo mais frequente em nível nacional ( 34,6 \%, correspondendo a 27 cursos ) e o conceito 2, o terceiro ( 12,8 \%, correspondendo a dez cursos ).
Houve, ainda, quatro cursos ( 5,1 \% ) que receberam conceito 5 e três cursos ( 3,8 \% ) que receberam conceito 1.
Dos 78 cursos de Geografia ( Bacharelado ), quatro ( 5,1 \% ) ficaram Sem Conceito ( SC ).
A região Norte participou com nove cursos ou 11,5 \% do total nacional.
Desses, seis cursos receberam os conceitos 2 e 3 ( três cursos cada ), o que \textbf{equivale} conjuntamente a 66,7 \% do total regional, além de serem os conceitos modais na região, como já comentado.
A dois cursos ( 22,2 \% ) atribuiu -se conceito 4.
Um curso ( 11,1 \% ) ficou Sem Conceito ( SC ) e nenhum curso recebeu o conceito 5 ou o conceito 1 nessa região.
A região Nordeste participou com 16 cursos ou 20,5 \% do total nacional.
Nessa região, sete cursos ( 43,8 \% em termos regionais ) obtiveram conceito 3, o conceito modal para a região, como já comentado.
Os conceitos 1 e 2 foram atribuídos, respectivamente, a um e a dois cursos ( 6,3 \% e 12,5 \%, respectivamente ).
Ainda nessa região, quatro cursos ( 25,0 \% ) foram avaliados com o conceito 4, e um curso ( 6,3 \% ), com conceito 5.
Um curso ficou Sem Conceito ( SC ).
Dos 27 cursos participantes da região Sudeste ( 34,6 \% do total nacional ), o conceito 4 foi o conceito modal, como já anteriormente assinalado, tendo sido obtido por 14 cursos ( 51,9 \% ).
O conceito 2 foi atribuído a um curso ( 3,7 \% ) e o conceito 3, a oito cursos ( 29,6 \% ).
Receberam o conceito 5, dois cursos ( 7,4 \% ).
Nessa região, nenhum curso recebeu conceito 1 e dois cursos ( 7,4 \% ) ficaram Sem Conceito ( SC ).
Os 17 cursos da região Sul corresponderam a 21,8 \% do total nacional.
Como já \textbf{indicado}, houve predominância do conceito 4, atribuído a 41,2 \% dos cursos da região ( conceito modal, equivalente a sete cursos ).
O conceito 1 foi atribuído a um curso ( 5,9 \% ), o conceito 2, a dois cursos ( 11,8 \% ), o conceito 3, a seis cursos ( 35,3 \% ) e o conceito 5, a um curso ( 5,9 \% ).
Nenhum curso ficou Sem Conceito ( SC ) na região Sul.
Já dos nove cursos participantes na região Centro-Oeste ( 11,5 \% do total nacional ), seis cursos receberam os conceitos 3 e 4 ( três cursos cada ), o que \textbf{equivale} conjuntamente a 66,7 \% do total regional, além de serem os conceitos modais na região, como já comentado.
A um curso ( 11,1 \% ) foi atribuído conceito 1, e a dois cursos ( 22,2 \% ), conceito 2.
O conceito \textbf{restante}, 5, não foi atribuído a nenhum curso.
Nessa região, nenhum curso ficou Sem Conceito ( SC ).
O Gráfico 5.1b apresenta a distribuição \textbf{cumulativa} do Conceito Enade dos cursos de Bacharelado segundo Grande Região.
Num gráfico de distribuição \textbf{cumulativa}, em cada valor de \textbf{abscissa} ( eixo x ), o valor apresentado no eixo das \textbf{ordenadas} ( eixo y ) é o valor acumulado da \textbf{variável}.
Por exemplo, a região Norte ( linha azul ) apresenta exatamente três quartos dos cursos com conceito 3 ou abaixo deste valor.
A região Nordeste ( linha verde ), por outro lado, apresenta exatamente dois terços dos cursos na mesma situação : com conceito 3 ou abaixo.
Nas regiões Norte ( linha azul ) e Centro-Oeste ( linha preta ), todos os cursos têm conceitos 4 ou abaixo disso, já que a \textbf{poligonal} \textbf{cumulativa} alcança 100 \% para o conceito 4.
Em linhas gerais, num gráfico de distribuição \textbf{cumulativa}, \textbf{poligonais} mais à esquerda representam Grandes Regiões com pior distribuição, e \textbf{poligonais} mais à direita, Grandes Regiões com melhores conceitos.
As regiões Norte ( linha azul ) e Centro-Oeste ( linha preta ) apresentam os cursos com conceitos mais baixos, ao passo que a região Sudeste ( linha vermelha ) apresenta a distribuição com valores maiores.
Cursos em IES situadas nas regiões Nordeste ( linha verde ) e Sul ( linha \textbf{roxa} ) apresentam uma situação intermediária. 5.2 CONCEITOS POR CATEGORIA ADMINISTRATIVA, POR MODALIDADE DE ENSINO E GRANDE REGIÃO A Tabela 5.2a apresenta a distribuição dos cursos de Licenciatura de Geografia participantes do Enade/2017, por Categoria Administrativa e Modalidade de ensino, de acordo com os conceitos por eles alcançados, segundo as Grandes Regiões brasileiras.
Dos 241 cursos participantes, 170 ( 70,5 \% ) eram \textbf{ministrados} em Instituições \textbf{Públicas} e 71 ( 29,5 \% ), em Instituições Privadas.
Quanto à Modalidade de ensino, 29 cursos ( 12,0 \% ) eram a Distância e 212 ( 88,0 \% ), Presenciais.
De acordo com as informações da Tabela 5.2a, em termos nacionais, entre Instituições \textbf{Públicas} e Privadas, observa -se que, dos 17 cursos avaliados com conceito 5, 14 eram oferecidos em IES Públicas e três, em IES Privadas.
Dos 170 cursos participantes de IES Públicas, o conceito 3 foi o valor modal, atribuído a 58 cursos ( 34,1 \% ).
Entre os demais cursos de Instituições \textbf{Públicas} participantes, sete cursos ( 4,1 \% ) foram avaliados com conceito 1, outros 43 cursos ( 25,3 \% ) obtiveram conceito 2, 47 cursos ( 27,6 \% ) receberam conceito 4 e, como já comentado, 14 cursos ( 8,2 \% ) receberam conceito 5.
Nesta categoria, somente um curso ( 0,6 \% ) ficou Sem Conceito ( SC ).
Na rede privada, o conceito modal foi 3, atribuído a 30 cursos, o correspondente a 42,3 \% dos 71 cursos da categoria.
Entre os demais cursos participantes, dois ( 2,8 \% ) receberam conceito 1, e 12 ( 16,9 \% ), conceito 2.
O conceito 4 foi atribuído a 19 cursos ( 26,8 \% ) e o conceito 5, a três cursos ( 4,2 \% ).
Nesta Categoria Administrativa, cinco cursos ( 7,0 \% ) ficaram Sem Conceito ( SC ).
Dos 212 cursos participantes da modalidade presencial, o conceito 3 foi o valor modal, atribuído a 75 cursos ( 35,4 \% ).
Entre os demais cursos nesta modalidade de ensino, 64 cursos ( 30,2 \% ) foram avaliados com conceito 4, 43 cursos ( 20,3 \% ) obtiveram conceito 2, 16 cursos ( 7,5 \% ) receberam conceito 5 e nove cursos ( 4,2 \% ) receberam conceito 1.
Nesta modalidade, 5 cursos ( 2,4 \% ) ficaram Sem Conceito ( SC ).
Entre os cursos a Distância, o conceito modal foi também 3, atribuído a 13 cursos, o correspondente a 44,8 \% dos 29 cursos da modalidade.
Entre os demais cursos participantes, dois ( 6,9 \% ) receberam conceito 4, um ( 3,4 \% ), conceito 5 e 12 ( 41,4 \% ), conceito 2.
Nessa Modalidade, um curso ( 3,4 \% ) ficou Sem Conceito ( SC ) e nenhum curso recebeu conceito 1.
Na análise por região, observa -se que, na região Norte, as Instituições Privadas participaram com três cursos ( 12,5 \% do total regional ), dos quais um ( 33,3 \% ) ficou Sem Conceito ( SC ) e os outros dois receberam conceito 3 ( 66,7 \% ).
As Instituições Públicas participaram com 21 cursos na região Norte ( 87,5 \% do total regional ), dos quais dez ( 47,6 \% ) obtiveram conceito 3, o valor modal.
O conceito 1 foi atribuído a dois cursos ( 9,5 \% ), o conceito 2 a sete cursos ( 33,3 \% ) e o conceito 4, a dois cursos ( 9,5 \% ).
Nenhum curso ficou Sem Conceito ( SC ) ou com conceito 5 nessa combinação de categoria e região.
Todos os cursos oferecidos na região Norte eram presenciais.
Na região Nordeste, a rede privada participou com apenas dois dos 63 cursos participantes, o equivalente a 3,2 \% do total da região, a menor proporção das Privadas entre as regiões.
Ambos os cursos oferecidos por IES Privadas no Nordeste foram avaliados com conceito 4.
As Instituições \textbf{Públicas} da região Nordeste participaram com 61 cursos, 96,8 \% do total da região, a maior proporção de cursos da rede pública entre as regiões.
Desses, 22 obtiveram conceito 3 ( 36,1 \% ), o valor modal, seguidos por outros 19 cursos, que receberam conceito 2 ( 31,1 \% ).
O conceito 4 foi atribuído a 14 cursos ( 23,0 \% ), e o conceito 1, a quatro cursos ( 6,6 \% ).
A dois cursos ( 3,3 \% ) foi atribuído o conceito 5.
Nenhum curso ficou Sem Conceito ( SC ).
A maioria dos cursos na região Nordeste foram presencias : 58 ( 92,1 \% ).
Para estes, a moda foi o conceito 3 com 20 cursos ( 34,5 \% ).
Os demais cursos presenciais da região foram alocados aos conceitos 1 ( quatro cursos, 6,9 \% ), 2 ( 16 cursos, 27,6 \% ), 4 ( também 16 cursos, 27,6 \% ) e 5 ( dois cursos, 3,4 \% ).
Nenhum curso presencial ficou Sem Conceito ( SC ).
Dos cinco cursos a Distância ( 7,9 \% ), três receberam conceito 2 ( 60,0 \% ), e dois receberam conceito 3 ( 40, \% ).
Nenhum curso a Distância da região foi alocado aos demais conceitos ou ficou Sem Conceito ( SC ).
Na região Sudeste, a proporção de cursos da rede privada, 53,6 \%, foi mais elevada do que a registrada nas demais regiões brasileiras, correspondendo a 45 ( 53,6 \% ) dos 84 cursos participantes na região.
Nessa categoria e região, o conceito modal foi 3, atribuído a 21 cursos ( 46,7 \% ).
Três cursos ( 6,7 \% ) ficaram Sem Conceito ( SC ).
Os demais foram avaliados com conceito 1 ( dois cursos, 4,4 \% ), conceito 2 ( seis cursos, 13,3 \% ), conceito 4 ( 11 cursos, 24,4 \% ) e conceito 5 ( dois cursos, 4,4 \% ).
Entre os 39 cursos ( 46,4 \% ) oferecidos por Instituições \textbf{Públicas} na região Sudeste, o conceito modal foi 4, atribuído a 16 cursos ( 41,0 \% ).
Os demais cursos receberam os conceitos 2 ( seis cursos, 15,4 \% ), 3 ( nove cursos, 23,1 \% ) e 5 ( oito cursos, 20,5 \% ).
Na região Sudeste, nesta categoria, nenhum curso ficou Sem Conceito ( SC ) ou recebeu o conceito 1.
Também na região Sudeste, a maioria dos cursos foram presencias : 70 ( 83,3 \% ).
Para estes, a moda foi o conceito 4 com 25 cursos ( 35,7 \% ).
Os demais cursos presenciais da região foram alocados aos conceitos 1 ( dois cursos, 2,9 \% ), 2 ( sete cursos, 10,0 \% ), 3 ( 24 cursos, 34,3 \% ) e 5 ( nove cursos, 12,9 \% ).
Três cursos presenciais ( 4,3 \% ) ficaram Sem Conceito ( SC ).
Para os 14 cursos a Distância ( 16,7 \% ), a moda também foi o conceito 3 ( 42,9 \% ).
Os demais cursos a Distância da região foram alocados aos conceitos 2 ( cinco cursos, 35,7 \% ), 4 ( dois cursos, 14,3 \% ).
Nenhum curso a Distância da região foi alocado ao conceito 1 ou ficou Sem Conceito ( SC ).
As Instituições Privadas participaram com 15 dos 38 cursos da região Sul, 39,5 \% do total regional.
Desses, sete ( 46,7 \% ) obtiveram conceito 3, o conceito modal.
Nesta combinação de Categoria Administrativa e Grande Região, os conceitos 2 e 4 foram atribuídos a três cursos ( 20,0 \% ) cada.
Um curso ( 6,7 \% ) recebeu conceito 5, e outro ( 6,7 \% ) ficou Sem Conceito ( SC ).
As Instituições \textbf{Públicas} na região Sul participaram com 23 cursos, 60,5 \% dos cursos da região.
O conceito modal foi 4, atribuído a 12 cursos ( 52,2 \% ).
Aos demais cursos foram atribuídos os conceitos 3 ( sete cursos, 30,4 \% ) e 5 ( três cursos, 13,0 \% ).
Um curso ( 4,3 \% ) da região Sul oferecido por IES Pública ficou Sem Conceito ( SC ), e nenhum recebeu os conceitos 1 ou 2.
Também na região Sul, a maioria dos cursos foram presenciais, correspondendo a 31 ( 81,6 \% ), por oposição aos sete cursos a Distância ( 18,4 \% ).
Para os cursos presenciais, a moda foi o conceito 4 com 15 cursos ( 48,4 \% ).
Os demais cursos presenciais da região foram alocados aos conceitos 2 ( um curso, 3,2 \% ), 3 ( dez cursos, 32,3 \% ) e 5 ( quatro cursos, 12,9 \% ).
Nenhum curso presencial foi alocado no conceito 1, e um curso ( 3,2 \% ) ficou Sem Conceito ( SC ).
Para os sete cursos a Distância, a moda foi o conceito 3 ( 57,1 \% ).
Outros dois cursos a Distância da região foram alocados ao conceito 2 ( 28,6 \% ).
Nenhum curso a Distância da região foi alocado aos demais conceitos, e um curso ( 14,3 \% ) ficou Sem Conceito ( SC ).
Na região Centro-Oeste, 6 dos 32 cursos participantes eram de Instituições Privadas ( 18,8 \% em termos regionais ).
Desses, três receberam conceito 2 e três receberam conceito 4.
Nenhum curso ficou Sem Conceito ( SC ) ou com conceitos 1, 3 ou 5.
Quanto aos 26 cursos oferecidos por Instituições \textbf{Públicas} na região Centro-Oeste ( 81,3 \% do total regional ), o conceito modal foi o conceito 2, com 11 cursos ( 42,3 \% ).
Os demais foram alocados aos conceitos 1 ( um curso, 3,8 \% ), 3 ( dez cursos, 38,5 \% ), 4 ( três cursos, 11,5 \% ) e 5 ( um curso, 3,8 \% ).
Nenhum curso ficou Sem Conceito ( SC ).
Nesta região, 29 cursos foram presenciais ( 90,6 \% ) e três, a Distância ( 9,4 \% ).
Para os cursos presenciais, a moda foi o conceito 2 com 12 cursos ( 41,4 \% ).
Os demais foram alocados aos conceitos 1 ( um curso, 3,4 \% ), 3 ( nove cursos, 31,0 \% ), 4 ( seis cursos, 20,7 \% ) e 5 ( um curso, 3,4 \% ).
Nenhum curso presencial ficou Sem Conceito ( SC ).
Para os três cursos a Distância ( 9,4 \% ), dois foram alocados ao conceito 2 ( 66,7 \% ), e um ao conceito 3 ( 33,3 \% ).
O Gráfico 5.2a apresenta a distribuição do Conceito Enade dos cursos de Licenciatura segundo a Categoria Administrativa da IES.
As duas \textbf{poligonais} se \textbf{cruzam}, com os cursos em IES Públicas ( linha azul ) apresentando uma \textbf{poligonal} menos inclinada do que aqueles em IES Privadas ( linha verde ) e, consequentemente, uma distribuição de Conceitos Enade com valores mais dispersos.
A Tabela 5.2b apresenta a distribuição dos cursos de Bacharelado participantes do Enade/2017 de Geografia, por Categoria Administrativa, de acordo com os conceitos por eles alcançados, segundo as Grandes Regiões brasileiras.
Dos 78 cursos participantes, 67 ( 85,9 \% ) eram \textbf{ministrados} em Instituições \textbf{Públicas} e 11 ( 14,1 \% ), em Instituições Privadas.
Quanto a Modalidade de ensino, 1 curso era a Distância e 77, presenciais.
De acordo ainda com as informações da Tabela 5.2b, em termos nacionais, entre Instituições \textbf{Públicas} e Privadas, observa -se que, dos quatro cursos avaliados com conceito 5, três eram oferecidos em IES Públicas e um, em IES Privada.
Dos 67 cursos participantes de IES Públicas, o conceito 4 foi o valor modal, atribuído a 27 cursos ( 40,3 \% ).
Entre os demais cursos de Instituições \textbf{Públicas} participantes, três cursos ( 4,5 \% ) foram avaliados com conceito 1, nove cursos ( 13,4 \% ) obtiveram conceito 2, 24 cursos ( 35,8 \% ) receberam conceito 3 e, como já comentado, três cursos ( 4,5 \% ) receberam conceito 5.
Nesta categoria, um curso ( 1,5 \% ) ficou Sem Conceito ( SC ).
Na rede privada, a moda ficou dividida entre os conceitos contíguos 3 e 4, com três cursos ( 27,3 \% ) cada.
Entre os demais cursos participantes, um ( 9,1 \% ) recebeu conceito 2, e outro ( 9,1 \% ), conceito 5.
Nenhum curso foi alocado ao conceito 1.
Nesta Categoria Administrativa, três cursos ( 27,3 \% ) ficaram Sem Conceito ( SC ).
Dos 77 cursos participantes da modalidade presencial ( 98,7 \% do total regional ), o conceito 4 foi o valor modal, atribuído a 29 cursos ( 37,7 \% ).
Entre os demais cursos nesta modalidade de ensino, 27 cursos ( 35,1 \% ) foram avaliados com conceito 3, dez cursos ( 13,0 \% ) obtiveram conceito 2, quatro cursos ( 5,2 \% ) receberam conceito 5 e três cursos ( 3,9 \% ) receberam conceito 1.
Nesta modalidade quatro cursos ( 5,2 \% ) ficaram Sem Conceito ( SC ).
Somente um curso a Distância foi avaliado para Geografia ( Bacharelado ) e este recebeu o conceito 4.
Este curso estava localizado na região Sudeste.
Na análise por região, observa -se que, na região Norte, não houve participação de cursos em Instituições Privadas.
Todos os nove cursos participantes foram de Instituições \textbf{Públicas}.
A moda ficou dividida em dois conceitos contíguos, 2 e 3, com três cursos ( 33,3 \% ) cada.
O conceito 4 foi atribuído a dois cursos ( 22,2 \% ).
Um curso ficou Sem Conceito ( SC ) e nenhum recebeu conceitos 1 ou 5 nessa combinação de categoria e região.
Todos os cursos oferecidos na região Norte eram presenciais.
Na região Nordeste, a rede pública concentrou 15 dos 16 cursos participantes, o equivalente a 93,8 \% do total da região.
O único curso oferecido por IES Privada no Nordeste, ficou Sem Conceito ( SC ).
Dos 15 cursos oferecidos por Instituições \textbf{Públicas} da região Nordeste, sete ( 46,7 \% ) obtiveram conceito 3, o valor modal, seguidos por outros quatro cursos ( 26,7 \% ), que receberam conceito 4 e dois cursos ( 13,3 \% ) com conceito 2.
Os conceitos 1 e 5 foram atribuídos a um curso ( 6,7 \% ) cada.
Nenhum curso ficou Sem Conceito ( SC ) nessa região.
Todos os cursos oferecidos na região Nordeste eram presenciais.
Na região Sudeste, a proporção de cursos da rede privada foi 29,6 \%, correspondendo a oito dos 27 cursos participantes.
Nessa categoria e região, o conceito modal foi 4, atribuído a três cursos ( 37,5 \% ).
Dois cursos ( 25,0 \% ) ficaram Sem Conceito ( SC ).
Os três cursos \textbf{restantes} foram avaliados com conceitos 2, 3 e 5 ( um curso, 12,5 \%, cada ).
Entre os 19 cursos ( 70,4 \% ) oferecidos por Instituições \textbf{Públicas} na região Sudeste, o conceito modal também foi 4, atribuído a 11 cursos ( 57,9 \% ).
Os demais cursos receberam os conceitos 3 ( sete cursos, 36,8 \% ) e 5 ( um curso, 5,3 \% ).
Na região Sudeste, nesta categoria, nenhum curso ficou Sem Conceito ( SC ).
Na região Sudeste, a maioria dos cursos foram presencias : 26 ( 96,3 \% ).
Para estes, a moda foi o conceito 4 com 13 cursos ( 50,0 \% ).
Os demais cursos presenciais da região foram alocados aos conceitos 2 ( um curso, 3,8 \% ), 3 ( oito cursos, 30,8 \% ) e 5 ( dois cursos, 7,7 \% ).
Nenhum curso presencial foi alocado ao conceito 1 e dois cursos ( 7,7 \% ) ficaram sem ( SC ).
O único curso a Distância foi alocado ao conceito 4, como já mencionado.
As Instituições Privadas participaram com 2 dos 17 cursos da região Sul, 11,8 \% do total regional.
Ambos foram alocados ao conceito 3.
As Instituições \textbf{Públicas} na região Sul participaram com 15 cursos, 88,2 \% dos cursos da região.
O conceito modal foi 4, atribuído a sete cursos ( 46,7 \% ).
Aos demais cursos foram atribuídos os conceitos 1 ( um curso, 6,7 \% ), 2 ( dois cursos, 13,3 \% ), 3 ( quatro cursos, 26,7 \% ) e 5 ( um curso, 6,7 \% ).
Nenhum curso da região Sul oferecido por IES Pública ficou Sem Conceito ( SC ).
Todos os cursos oferecidos na região Sul eram presenciais.
Na região Centro-Oeste, todos os nove cursos avaliados eram de Instituições \textbf{Públicas} e Presenciais.
A moda ficou dividida entre dois conceitos contíguos : 3 e 4, com três cursos ( 33,3 \% ) cada.
Um curso ( 11,1 \% ) apresentou conceito 1 e dois ( 22,2 \% ) ficaram com conceitos 2.
Nenhum curso ficou Sem Conceito ( SC ) ou com conceito 5.
O Gráfico 5.2b apresenta a distribuição do Conceito Enade dos cursos de Bacharelado segundo a Categoria Administrativa da IES.
Os cursos em IES Públicas ( linha azul ) apresentam uma \textbf{poligonal} mais à esquerda do que os em IES Privadas ( linha verde ) e, consequentemente, uma distribuição de Conceitos Enade com valores menores. 5.3 CONCEITOS POR ORGANIZAÇÃO ACADÊMICA E GRANDE REGIÃO Na Tabela 5.3a, encontra -se a distribuição dos conceitos atribuídos aos cursos de Licenciatura participantes do Enade/2017 na Área de Geografia, por Organização Acadêmica, segundo as Grandes Regiões brasileiras.
Dos 241 cursos participantes, 190 eram oferecidos em Universidades, 19 em Centros Universitários, 22 em Faculdades e 10 em CEFET/IFET.
Esta distribuição corresponde a, respectivamente, 78,8 \%, 7,9 \%, 9,1 \% e 4,1 \% dos cursos.
De acordo com os dados apresentados, de todos os 17 cursos avaliados com conceito 5, 14 eram vinculados a Universidades.
Os cursos vinculados a esse tipo de Organização Acadêmica tiveram o conceito 3 como conceito modal, atribuído a 69 cursos ( 36,3 \% ).
Os demais cursos vinculados a Universidades receberam os conceitos 1 ( seis cursos, 3,2 \% ), 2 ( 43 cursos, 22,6 \% ), 4 ( 56 cursos, 29,5 \% ) e conceito 5 ( 14 cursos, 7,4 \%, como já mencionado ).
Ficaram Sem Conceito ( SC ), dois cursos ( 1,1 \% ).
Entre os cursos em Centros Universitários, o conceito modal foi também 3, atribuído a seis cursos ( 31,6 \% ).
Quatro cursos ( 21,1 \% ) vinculado a esse tipo de Organização Acadêmica ficaram Sem Conceito ( SC ), e os outros cursos receberam os conceitos 2 ( quatro cursos, 21,1 \% ) e 4 ( cinco cursos, 26,3 \% ).
Nenhum curso recebeu os conceitos 1 e 5.
Dos 22 cursos mantidos por Faculdades, dez ( 45,5 \% ) receberam conceito 3, o conceito modal.
O segundo conceito mais frequentemente atribuído aos cursos oferecidos por essa categoria de IES foi o conceito 2, recebido por seis cursos ( 27,3 \% ).
Na sequência, três cursos obtiveram conceito 1 ( 13,6 \% ), dois cursos receberam conceito 4 ( 9,1 \% ) e um curso, conceito 5 ( 4,5 \% ).
Nenhum curso ficou Sem Conceito ( SC ).
Dos 10 cursos mantidos por CEFET/IFET, seis ( 60,0 \% ) receberam conceitos 3 e 4 ( três em cada ), os conceitos modais.
O segundo conceito mais frequentemente atribuído aos cursos oferecidos por essa categoria de IES foi o conceito 2, recebido por dois cursos ( 20,0 \% ).
Outros dois cursos receberam conceito 5 ( 20,0 \% ).
Nenhum curso ficou Sem Conceito ( SC ) ou recebeu o conceito 1.
Considerando -se separadamente as regiões brasileiras, verifica -se que, na região Norte, as Universidades participaram com um 21 ( 87,5 \% ) dos 24 cursos participantes.
O conceito modal entre os cursos de Universidades, 3, foi alocado a 11 cursos ( 52,4 \% ).
Os demais cursos receberam os conceitos 1 ( dois cursos, 9,5 \% ), 2 ( sete cursos, 33,3 \% ) e 4 ( um curso, 4,8 \% ).
Nenhum curso ficou Sem Conceito ( SC ) ou recebeu conceito 5 nesta combinação de categoria e região.
Os Centros Universitários da região Norte foram representados por somente um curso que ficou Sem Conceito ( SC ).
As Faculdades participaram também com um curso, que recebeu conceito 3.
Os CEFET/IFET participaram igualmente com um curso, e que recebeu conceito 4.
Na região Nordeste, as Universidades participaram com 54 ( 85,7 \% ) dos 63 cursos da Área de Geografia ( Licenciatura ) oferecidos na região, dos quais nenhum ficou Sem Conceito ( SC ).
Nota -se que tais cursos apresentaram uma distribuição com moda no conceito 3, atribuído a 21 cursos ( 38,9 \% ).
Obtiveram os conceitos 2 e 4, respectivamente, 16 ( 29,6 \% ) e 14 ( 25,9 \% ) cursos.
Dois cursos ( 3,7 \% ) receberam o conceito 1 e um curso, o conceito 5 ( 1,9 \% ).
Os Centros Universitários \textbf{contaram} com um curso participante na região Nordeste.
Esse obteve conceito 4.
As Faculdades foram representadas por três cursos na região que apresentaram moda no conceito 1, com dois cursos alocados.
O conceito 2 foi atribuído ao outro curso da região.
Os CEFET/IFET participaram com cinco cursos, os quais receberam conceitos 2 ( dois cursos ), 3, 4 e 5 ( um curso cada ).
Na região Sudeste, as Universidades concentraram 55 ( 65,5 \% ) dos 84 cursos de Geografia ( Licenciatura ) da região, a menor proporção dessa categoria entre as regiões.
Entre os cursos em Universidades, nessa região, o conceito modal foi 4, atribuído a 23 cursos ( 41,8 \% ).
Os demais cursos receberam os conceitos 1 ( um curso, 1,8 \% ), 2 ( oito cursos, 14,5 \% ), 3 ( 14 cursos, 25,5 \% ) e 5 ( oito cursos, 14,5 \% ).
Um curso ( 1,8 \% ) ficou Sem Conceito ( SC ).
Os Centros Universitários participaram com 11 cursos ( 13,1 \% ) na região Sudeste, dos quais cinco obtiveram o conceito modal 3 ( 45,5 \% ).
O conceito 2 foi atribuído a dois cursos ( 18,2 \% ) e o conceito 4, também a dois cursos ( 18,2 \% ).
Outros dois cursos ( 18,2 \% ) ficaram Sem Conceito ( SC ) nessa combinação de organização acadêmica e região.
Nenhum curso foi alocado aos conceitos 1 e 5.
Já as Faculdades foram representadas por 14 cursos na região Sudeste, 16,7 \% da representação regional, a maior para essa categoria entre as regiões, dos quais nenhum ficou Sem Conceito ( SC ).
Aos demais cursos foram atribuídos os conceitos 1 ( um curso, 7,1 \% ), 2 ( dois cursos, 14,3 \% ), 3 ( nove cursos, 64,3 \%, o conceito modal ), 4 ( um curso, 7,1 \% ) e 5 ( um curso, 7,1 \% ).
Os CEFET/IFET participaram com quatro cursos, os quais receberam conceitos 3 ( dois cursos, o conceito modal ), 4 e 5 ( um curso cada ).
Dos 38 cursos da região Sul, 32 ( 84,2 \% ) eram de Universidades, para os quais o conceito modal foi 4, atribuído a 14 cursos ( 43,8 \% ).
Nesse tipo de organização, um curso ( 3,1 \% ) ficou Sem Conceito ( SC ) e nenhum curso recebeu conceito 1 ou 2.
Os demais cursos receberam os conceitos 3 ( 13 cursos, 40,6 \% ) e 5 ( quatro cursos, 12,5 \% ).
Dos cinco cursos participantes de Centros Universitários da região Sul, que representaram 13,2 \% dos cursos da região, a maior proporção dessa categoria entre as regiões, dois cursos ficaram com conceito 2, um curso obteve conceito 3, um curso conceito 4 e um curso ficou Sem Conceito ( SC ).
O único curso vinculado a Faculdades na região Sul recebeu conceito 2.
Nesta região não foram avaliados cursos em CEFET/IFET.
Na região Centro-Oeste, 28 ( 87,5 \% ) dos 32 cursos eram de Universidades.
Nesse tipo de organização, nenhum curso ficou Sem Conceito ( SC ).
Os conceitos 1 e 5 foram atribuídos a um curso ( 3,6 \% ) cada, o conceito 2, o conceito modal, a 12 cursos ( 42,9 \% ), o conceito 3, a dez cursos ( 35,7 \% ) e o conceito 4 a quatro cursos ( 14,3 \% ).
Os Centros Universitários da região Centro-Oeste \textbf{contaram} com somente um curso, que obteve conceito 4.
Dos três cursos oferecidos em Faculdades na região Centro-Oeste, dois receberam o conceito modal 2, e um curso ficou com conceito 4.
Nesta região não foram avaliados cursos em CEFET/IFET.
O Gráfico 5.3a apresenta a distribuição do Conceito Enade dos cursos de Licenciatura segundo a Organização Acadêmica da IES.
Os cursos em CEFET/IFET ( linha \textbf{roxa} ) apresentam uma \textbf{poligonal} mais à direita do que as demais, \textbf{denotando} uma melhor distribuição de Conceitos Enade.
No outro extremo, temos os cursos oferecidos em Faculdades ( linha vermelha ) com a \textbf{poligonal} mais à esquerda e, consequentemente, uma distribuição de Conceitos Enade com menores valores.
IES de Centros Universitários ( linha verde ) e em Universidades ( linha azul ) apresentam uma situação intermediária.
Na Tabela 5.3b, encontra -se a distribuição dos conceitos atribuídos aos cursos de Bacharelado participantes do Enade/2017 na Área de Geografia, por Organização Acadêmica, segundo as Grandes Regiões brasileiras.
Dos 78 cursos participantes, 76 eram oferecidos em Universidades e dois em Centros Universitários.
Não houve cursos avaliados na área de Geografia ( Bacharelado ) em Faculdades, ou CEFET/IFET.
Esta distribuição corresponde a, respectivamente, 97,4 \% e 2,6 \% dos cursos.
De acordo com os dados apresentados, os quatro cursos avaliados com conceito 5 eram vinculados a Universidades.
Os cursos vinculados a esse tipo de Organização Acadêmica tiveram o conceito 4 como conceito modal, atribuído a 30 cursos ( 39,5 \% ).
Os demais cursos vinculados a Universidades receberam os conceitos 1 ( três cursos, 3,9 \% ), 2 ( nove cursos, 11,8 \% ), 3 ( 26 cursos, 34,2 \% ) e conceito 5 ( quatro cursos, 5,3 \%, como já mencionado ).
Ficaram Sem Conceito ( SC ), quatro cursos ( 5,3 \% ).
Os dois cursos em Centros Universitários receberam conceitos 2 e 3.
Considerando -se separadamente as regiões brasileiras, verifica -se que, na região Norte, as Universidades participaram com a totalidade dos cursos.
Seis cursos receberam os conceitos modais 2 e 3 ( três cursos cada ).
Os demais cursos receberam o conceito 4 ( dois cursos ) e um curso ficou Sem Conceito ( SC ).
Os conceitos 1 e 5 não foram atribuídos a nenhum curso.
Não houve participação de cursos nem em Centros Universitários, nem em Faculdades e nem em CEFET/IFET na região Norte.
Na região Nordeste, todos os cursos também eram oferecidos por Universidades, que participaram com os 16 cursos da Área de Geografia ( Bacharelado ) oferecidos na região, um dos quais ( 6,3 \% ) ficou Sem Conceito ( SC ).
A moda foi o conceito 3 com sete cursos ( 43,8 \% ).
Obtiveram os conceitos 2 e 4, respectivamente, dois ( 12,5 \% ) e quatro ( 25,0 \% ) cursos.
Os dois cursos \textbf{restantes} foram alocados aos conceitos 1 e 5 ( 6,3 \% para cada conceito ).
Não houve participação de cursos nem em Centros Universitários, nem em Faculdades e nem em CEFET/IFET na região Nordeste.
Na região Sudeste, as Universidades concentraram 25 ( 92,6 \% ) dos 27 cursos de Geografia ( Bacharelado ) da região.
Entre os cursos em Universidades, nessa região, o conceito modal foi 4, atribuído a 14 cursos ( 56,0 \% ).
Os demais cursos receberam os conceitos 3 ( sete cursos, 28,0 \% ) e 5 ( dois cursos, 8,0 \% ).
Dois cursos ( 8,0 \% ) ficaram Sem Conceito ( SC ).
Nenhum curso recebeu os conceitos 1 e 2.
Os Centros Universitários participaram com dois cursos na região Sudeste, dos quais um obteve conceito 2 e o outro, conceito 3.
Não houve participação de cursos em Centros Universitários, Faculdades ou CEFET/IFET nesta região.
Todos os 17 cursos da região Sul estavam concentrados exclusivamente em Universidades, para os quais o conceito modal foi 4, atribuído a sete cursos ( 41,2 \% ).
Nesse tipo de organização, nenhum curso ficou Sem Conceito ( SC ).
Os demais cursos receberam os conceitos 1 ( um curso, 5,9 \% ), 2 ( dois cursos, 11,8 \% ), 3 ( seis cursos, 35,3 \% ) e 5 ( um curso, 5,9 \% ).
Na região Centro-Oeste, também todos os nove cursos eram de Universidades.
Nesse tipo de organização, nenhum curso ficou Sem Conceito ( SC ) ou alcançou o conceito máximo 5.
Os conceitos 3 e 4 ( conceitos modais ) foram atribuídos a três cursos ( 33,3 \% ) cada, o conceito 2, a dois cursos ( 22,2 \% ) e o conceito 1, a um curso ( 11,1 \% ).
Nenhum curso alcançou o conceito máximo, 5.
Como já comentado, não houve participação de cursos nem em Centros Universitários, nem em Faculdades e nem em CEFET/IFET na região Centro-Oeste.
O Gráfico 5.3b apresenta a distribuição do Conceito Enade dos cursos de Bacharelado segundo a Organização Acadêmica da IES.
Os cursos em Universidades ( linha azul ) apresentam uma \textbf{poligonal} mais à direita do que a de Centros Universitários ( linha verde ), \textbf{denotando} uma melhor distribuição de notas nos cursos oferecidos em Universidades ( linha azul ).

% matched lemmas: abscissa, contar, cruzar, cumulativo, denotar, equivaler, indicado, intermédio, ministrar, ordenado, panorama, poligonal, pública, restante, roxa, tabular, variável
\end{textsample}
